% GENERATED FILE. Edit canonical lessons, then run npm run generate:books.
% canonical-source-hash: fnv1a64:fb3d0337877dd8ef
% canonical-lessons: TE-C376-spandana, TE-C376-prarambham, TE-C376-mugimpu, TE-C376-ghattam, TE-C376-yadrcchikanga

\chapter{How It Began and Ended}
\label{ch:te-how-it-began-and-ended}
\hlchaptermodality{376}

\begin{chapteropening}
\textbf{By the end of this chapter:} \emph{I can say a response, a start, an ending, an episode and by chance.}

Complete the last lesson of chapter 376: I can say a response, a start, an ending, an episode and by chance.
\end{chapteropening}

\section[\texorpdfstring{\te{స్పందన} (spandana)}{స్పందన (spandana)}]{\te{స్పందన} (spandana) — a response; a reaction}

\label{lesson:TE-C376-spandana}



\begin{quote}
Before the new one: say the Telugu for to submit a report, then the Telugu for to respond.
\end{quote}

\subsection*{You'll want to know: \te{స్పందన}}
\textbf{\te{స్పందన}} — \emph{spandana} — \textquotedblleft{}a response; a reaction\textquotedblright{}.

\begin{quote}
\te{ఆ} \te{పుస్తకానికి} \te{మంచి} \te{స్పందన} \te{వచ్చింది}.
\end{quote}

\begin{quote}
\emph{ā pustakāniki mañci spandana vaccindi.}
\end{quote}

\begin{quote}
That book got a good response.
\end{quote}

The noun of \te{స్పందించు}.

\begin{cousinweb}
\textbf{\te{స్పందన}} (\emph{spandana}) is borrowed from Sanskrit \emph{spandana} (Monier-Williams), from \emph{spand}, \textquotedblleft{}to quiver, throb\textquotedblright{}.
\end{cousinweb}

\begin{grammarlens}[title={What you've built}]
One more word for telling what happened.
\end{grammarlens}

\subsection*{Your turn}
\begin{itemize}
\raggedright
  \item \emph{Say it:} \emph{spandana}
  \item \emph{Say it:} \emph{spandana}, once more
  \item \emph{Say it:} \emph{ā pustakāniki mañci spandana vaccindi.} — That book got a good response.
  \item \emph{Recall:} say the Telugu for to submit a report, then the Telugu for to respond, then say \emph{spandana} again
\end{itemize}

\begin{tcolorbox}[breakable,colback=teal!4,colframe=teal!35!black,title={Before you move on}]
What does \te{స్పందన} mean? (\textquotedblleft{}A response; a reaction\textquotedblright{}.) Say it once more.
\end{tcolorbox}
\section[\texorpdfstring{\te{ప్రారంభం} (prārambhaṁ)}{ప్రారంభం (prārambhaṁ)}]{\te{ప్రారంభం} (prārambhaṁ) — a start; a beginning}

\label{lesson:TE-C376-prarambham}



\begin{quote}
Before the new one: say the Telugu for to respond, then the Telugu for a response.
\end{quote}

\subsection*{You'll want to know: \te{ప్రారంభం}}
\textbf{\te{ప్రారంభం}} — \emph{prārambhaṁ} — \textquotedblleft{}a start; a beginning\textquotedblright{}.

\begin{quote}
\te{కథ} \te{ప్రారంభం} \te{చాలా} \te{బాగుంది}.
\end{quote}

\begin{quote}
\emph{katha prārambhaṁ cālā bāgundi.}
\end{quote}

\begin{quote}
The start of the story is very good.
\end{quote}

A formal word; in speech \te{మొదలు}.

\begin{cousinweb}
\textbf{\te{ప్రారంభం}} (\emph{prārambhaṁ}) is borrowed from Sanskrit \emph{prārambha}, \textquotedblleft{}commencement, beginning\textquotedblright{} (Monier-Williams).
\end{cousinweb}

\begin{grammarlens}[title={What you've built}]
One more word for telling what happened.
\end{grammarlens}

\subsection*{Your turn}
\begin{itemize}
\raggedright
  \item \emph{Say it:} \emph{prārambhaṁ}
  \item \emph{Say it:} \emph{prārambhaṁ}, once more
  \item \emph{Say it:} \emph{katha prārambhaṁ cālā bāgundi.} — The start of the story is very good.
  \item \emph{Recall:} say the Telugu for to respond, then the Telugu for a response, then say \emph{prārambhaṁ} again
\end{itemize}

\begin{tcolorbox}[breakable,colback=teal!4,colframe=teal!35!black,title={Before you move on}]
What does \te{ప్రారంభం} mean? (\textquotedblleft{}A start; a beginning\textquotedblright{}.) Say it once more.
\end{tcolorbox}
\section[\texorpdfstring{\te{ముగింపు} (mugimpu)}{ముగింపు (mugimpu)}]{\te{ముగింపు} (mugimpu) — an ending; a conclusion}

\label{lesson:TE-C376-mugimpu}



\begin{quote}
Before the new one: say the Telugu for a response, then the Telugu for a start.
\end{quote}

\subsection*{You'll want to know: \te{ముగింపు}}
\textbf{\te{ముగింపు}} — \emph{mugimpu} — \textquotedblleft{}an ending; a conclusion\textquotedblright{}.

\begin{quote}
\te{కానీ} \te{కథ} \te{ముగింపు} \te{నాకు} \te{నచ్చలేదు}.
\end{quote}

\begin{quote}
\emph{kānī katha mugimpu nāku naccalēdu.}
\end{quote}

\begin{quote}
But I did not like the story's ending.
\end{quote}

The noun of \te{ముగించు}, to finish.

\begin{cousinweb}
\textbf{\te{ముగింపు}} (\emph{mugimpu}) is native Dravidian: DEDR 4891 lists Telugu \emph{mugimpu}, \textquotedblleft{}finishing, end\textquotedblright{}, with \emph{mugincu}, \textquotedblleft{}to finish\textquotedblright{}, and Tamil \emph{muki}, \textquotedblleft{}to end\textquotedblright{}.
\end{cousinweb}

\begin{grammarlens}[title={What you've built}]
One more word for telling what happened.
\end{grammarlens}

\subsection*{Your turn}
\begin{itemize}
\raggedright
  \item \emph{Say it:} \emph{mugimpu}
  \item \emph{Say it:} \emph{mugimpu}, once more
  \item \emph{Say it:} \emph{kānī katha mugimpu nāku naccalēdu.} — But I did not like the story's ending.
  \item \emph{Recall:} say the Telugu for a response, then the Telugu for a start, then say \emph{mugimpu} again
\end{itemize}

\begin{tcolorbox}[breakable,colback=teal!4,colframe=teal!35!black,title={Before you move on}]
What does \te{ముగింపు} mean? (\textquotedblleft{}An ending; a conclusion\textquotedblright{}.) Say it once more.
\end{tcolorbox}
\section[\texorpdfstring{\te{ఘట్టం} (ghaṭṭaṁ)}{ఘట్టం (ghaṭṭaṁ)}]{\te{ఘట్టం} (ghaṭṭaṁ) — an episode; a phase}

\label{lesson:TE-C376-ghattam}



\begin{quote}
Before the new one: say the Telugu for a start, then the Telugu for an ending.
\end{quote}

\subsection*{You'll want to know: \te{ఘట్టం}}
\textbf{\te{ఘట్టం}} — \emph{ghaṭṭaṁ} — \textquotedblleft{}an episode; a phase\textquotedblright{}.

\begin{quote}
\te{కథలో} \te{ఆ} \te{ఘట్టం} \te{నాకు} \te{బాగా} \te{గుర్తుంది}.
\end{quote}

\begin{quote}
\emph{kathalō ā ghaṭṭaṁ nāku bāgā gurtundi.}
\end{quote}

\begin{quote}
I remember that episode of the story well.
\end{quote}

Plural: \te{ఘట్టాలు}, \emph{ghaṭṭālu}.

Also a landing place; in a story, a stage of it.

\begin{cousinweb}
\textbf{\te{ఘట్టం}} (\emph{ghaṭṭaṁ}) is borrowed from Sanskrit \emph{ghaṭṭa}, \textquotedblleft{}a landing place, a quay\textquotedblright{} (CDIAL 4414); CDIAL gives its basic sense as a passage.
\end{cousinweb}

\begin{grammarlens}[title={What you've built}]
One more word for telling what happened.
\end{grammarlens}

\subsection*{Your turn}
\begin{itemize}
\raggedright
  \item \emph{Say it:} \emph{ghaṭṭaṁ}
  \item \emph{Say it:} \emph{ghaṭṭaṁ}, once more
  \item \emph{Say it:} \emph{kathalō ā ghaṭṭaṁ nāku bāgā gurtundi.} — I remember that episode of the story well.
  \item \emph{Recall:} say the Telugu for a start, then the Telugu for an ending, then say \emph{ghaṭṭaṁ} again
\end{itemize}

\begin{tcolorbox}[breakable,colback=teal!4,colframe=teal!35!black,title={Before you move on}]
What does \te{ఘట్టం} mean? (\textquotedblleft{}An episode; a phase\textquotedblright{}.) Say it once more.
\end{tcolorbox}
\section[\texorpdfstring{\te{యాదృచ్ఛికంగా} (yādṛcchikaṅgā)}{యాదృచ్ఛికంగా (yādṛcchikaṅgā)}]{\te{యాదృచ్ఛికంగా} (yādṛcchikaṅgā) — by chance}

\label{lesson:TE-C376-yadrcchikanga}



\begin{quote}
Before the new one: say the Telugu for an ending, then the Telugu for an episode.
\end{quote}

\subsection*{You'll want to know: \te{యాదృచ్ఛికంగా}}
\textbf{\te{యాదృచ్ఛికంగా}} — \emph{yādṛcchikaṅgā} — \textquotedblleft{}by chance\textquotedblright{}.

\begin{quote}
\te{మేము} \te{యాదృచ్ఛికంగా} \te{ఒకే} \te{రైలులో} \te{కలిశాము}.
\end{quote}

\begin{quote}
\emph{mēmu yādṛcchikaṅgā okē railulō kaliśāmu.}
\end{quote}

\begin{quote}
We met by chance on the same train.
\end{quote}

A bookish word; \te{అనుకోకుండా} says much the same in speech.

\begin{cousinweb}
\textbf{\te{యాదృచ్ఛికంగా}} (\emph{yādṛcchikaṅgā}) is built on Sanskrit \emph{yādṛcchika}, \textquotedblleft{}spontaneous, accidental\textquotedblright{} (Monier-Williams), with the Telugu adverb ending \emph{-gā}.
\end{cousinweb}

\begin{grammarlens}[title={What you've built}]
One more word for telling what happened.
\end{grammarlens}

\subsection*{Your turn}
\begin{itemize}
\raggedright
  \item \emph{Say it:} \emph{yādṛcchikaṅgā}
  \item \emph{Say it:} \emph{yādṛcchikaṅgā}, once more
  \item \emph{Say it:} \emph{mēmu yādṛcchikaṅgā okē railulō kaliśāmu.} — We met by chance on the same train.
  \item \emph{Recall:} say the Telugu for an ending, then the Telugu for an episode, then say \emph{yādṛcchikaṅgā} again
\end{itemize}

\begin{tcolorbox}[breakable,colback=teal!4,colframe=teal!35!black,title={Before you move on}]
What does \te{యాదృచ్ఛికంగా} mean? (\textquotedblleft{}By chance\textquotedblright{}.) Say it once more.
\end{tcolorbox}
